\subsection{\texorpdfstring{The homomorphism \(\operatorname{Ext}_{\mathbf{A}}(\mathbf{N}, \mathbf{P}) \otimes \operatorname{Ext}_{\mathbf{A}}(\mathbf{M}, \mathbf{N}) \to \operatorname{Ext}_{\mathbf{A}}(\mathbf{M}, \mathbf{P})\)}{The homomorphism \textbackslash operatorname\{Ext\}\_\{\textbackslash mathbf\{A\}\}(\textbackslash mathbf\{N\}, \textbackslash mathbf\{P\}) \textbackslash otimes \textbackslash operatorname\{Ext\}\_\{\textbackslash mathbf\{A\}\}(\textbackslash mathbf\{M\}, \textbackslash mathbf\{N\}) \textbackslash to \textbackslash operatorname\{Ext\}\_\{\textbackslash mathbf\{A\}\}(\textbackslash mathbf\{M\}, \textbackslash mathbf\{P\})}}

Let M and N be two left A-modules. Consider the homologisms \(p_{\mathbf{M}}: \mathbf{L}(\mathbf{M}) \to \mathbf{M}\) , \(e_{\mathbf{M}}: \mathbf{M} \to \mathbf{l}(\mathbf{M})\) and \(e_{\mathbf{M}} \circ p_{\mathbf{M}}: \mathbf{L}(\mathbf{M}) \to \mathbf{l}(\mathbf{M})\) ; by prop. 4 of X, p.~86, we deduce from this a bijective homomorphism

\[
a _ {\mathrm{M}, \mathrm{N}} = \mathrm{H} (\operatorname{Homgr} _ {\mathrm{A}} (e _ {\mathrm{M}} \circ p _ {\mathrm{M}}, 1)): \mathrm{H} (\operatorname{Homgr} _ {\mathrm{A}} (\mathrm{I} (\mathrm{M}), \mathrm{I} (\mathrm{N}))) \rightarrow \operatorname{Ext} _ {\mathrm{A}} (\mathrm{M}, \mathrm{N}).
\]

Recall moreover (X, p.~82) that \(\mathrm{H}^n (\mathrm{Homgr}_{\mathbf{A}}(\mathrm{I}(\mathbf{M}),\mathrm{I}(\mathbf{N})))\) is the set of homotopy classes of morphisms of ascending degree \(n\) from the complex I(M) into the complex I(N). For example, if \(f\in \mathrm{Hom}_{\mathbf{A}}(\mathbf{M},\mathbf{N})\) the homotopy class of I(f) is sent by \(a_{\mathbf{M},\mathbf{N}}\) to \(f\) .

If P is a third left A-module, we have, according to X, p.~99, a canonical k-homomorphism

\[
\mathrm{H} (\operatorname{Homgr} _ {\mathbf {A}} (I (N), I (P))) \otimes_ {k} \mathrm{H} (\operatorname{Homgr} _ {\mathbf {A}} (I (M), I (N))) \rightarrow \mathrm{H} (\operatorname{Homgr} _ {\mathbf {A}} (I (M), I (P)))
\]

from which we deduce by transport of structure via the isomorphisms \(a_{\mathbf{N},\mathbf{P}}, a_{\mathbf{M},\mathbf{N}}, a_{\mathbf{M},\mathbf{P}}\) a \(k\) -homomorphism (composition homomorphism):

\[
c _ {\mathbf {M}, \mathbf {N}, \mathbf {P}}: \operatorname{Ext} _ {\mathbf {A}} (\mathbf {N}, \mathbf {P}) \otimes_ {k} \operatorname{Ext} _ {\mathbf {A}} (\mathbf {M}, \mathbf {N}) \rightarrow \operatorname{Ext} _ {\mathbf {A}} (\mathbf {M}, \mathbf {P}).
\]

This corresponds to a k-bilinear map.

\[
\operatorname{Ext} _ {\mathbf {A}} (\mathbf {N}, \mathbf {P}) \times \operatorname{Ext} _ {\mathbf {A}} (\mathbf {M}, \mathbf {N}) \rightarrow \operatorname{Ext} _ {\mathbf {A}} (\mathbf {M}, \mathbf {P})\tag{1}
\]

which decomposes into homogeneous components

\[
\operatorname{Ext} _ {\mathbf {A}} ^ {i} (\mathbf {N}, \mathbf {P}) \times \operatorname{Ext} _ {\mathbf {A}} ^ {j} (\mathbf {M}, \mathbf {N}) \rightarrow \operatorname{Ext} _ {\mathbf {A}} ^ {i + j} (\mathbf {M}, \mathbf {P}).\tag{2}
\]

If \(u \in \text{Ext}_A(\mathbf{N}, \mathbf{P})\) , \(v \in \text{Ext}_A(\mathbf{M}, \mathbf{N})\) , the image of \((u, v)\) under (1) is called the composition product of \(u\) and \(v\) and is denoted \(u \circ v\) . If \(g\) (resp. \(f\) ) is a morphism of complexes of ascending degree \(j\) (resp. \(i\) ) from I(M) to I(N) (resp. from I(N) to I(P)), of homotopy class \(\overline{g}\) (resp. \(\overline{f}\) ), then the composition product \(a_{\mathbf{N},\mathbf{P}}(\overline{f}) \circ a_{\mathbf{M},\mathbf{N}}(\overline{g})\) is the image under \(a_{\mathbf{M},\mathbf{P}}\) of the homotopy class of the morphism \(f \circ g\) from I(M) into I(P).

Example 1. --- If \(u \in \operatorname{Hom}_{\mathbf{A}}(\mathbf{N}, \mathbf{P})\) , \(v \in \operatorname{Hom}_{\mathbf{A}}(\mathbf{M}, \mathbf{N})\) , \(u \circ v\) is the composite of \(u\) and of \(v\) .

Example 2. --- If \(u \in \mathrm{Hom}_{\mathbf{A}}(\mathbf{N}, \mathbf{P})\) , \(v \in \mathrm{Ext}_{\mathbf{A}}(\mathbf{M}, \mathbf{N})\) , then

\[
u \circ v = \operatorname{Ext} _ {A} (1 _ {M}, u) (v) \in \operatorname{Ext} _ {A} (M, P);
\]

likewise, if \(u \in \operatorname{Ext}_{\mathbf{A}}(\mathbf{N}, \mathbf{P})\) , \(v \in \operatorname{Hom}_{\mathbf{A}}(\mathbf{M}, \mathbf{N})\) , then

\[
u \circ v = \operatorname{Ext} _ {A} (v, 1 _ {p}) (u) \in \operatorname{Ext} _ {A} (M, P).
\]

This follows from the definitions and remarks in X, p.~88.

If Q, M, N, P are four left A-modules, and if

\[
u \in \operatorname{Ext} _ {\mathbf {A}} (\mathbf {N}, \mathrm{P}), \quad v \in \operatorname{Ext} _ {\mathbf {A}} (\mathbf {M}, \mathbf {N}), \quad w \in \operatorname{Ext} _ {\mathbf {A}} (\mathbf {Q}, \mathbf {M}),
\]

then \((u \circ v) \circ w = u \circ (v \circ w)\) : the composition product is associative; we will therefore write the composites of several elements without parentheses. In particular, according to Example 2:

Example 3. — Let M, N, M', N' be four left A-modules. If \(u \in \operatorname{Ext}_{\mathbf{A}}(\mathbf{M}, \mathbf{N})\) , \(f \in \operatorname{Hom}_{\mathbf{A}}(\mathbf{M}', \mathbf{M})\) , \(g \in \operatorname{Hom}_{\mathbf{A}}(\mathbf{N}, \mathbf{N}')\) , then

\[
\operatorname{Ext} _ {\mathbf {A}} (f, g) (u) = g \circ u \circ f \in \operatorname{Ext} _ {\mathbf {A}} (\mathbf {M} ^ {\prime}, \mathbf {N} ^ {\prime}).\tag{3}
\]

This gives a new proof of the \(k\) -bilinearity of the mapping \((f, g) \to \operatorname{Ext}_{\mathbf{A}}(f, g)\) (X, p.~88, prop. 6).

